\documentclass[11pt]{article}

\usepackage[T1]{fontenc}
\usepackage[utf8]{inputenc}
\usepackage[english]{babel}
\usepackage{amsmath}
\usepackage{amssymb}
\usepackage{graphicx}
\usepackage{geometry}
\geometry{margin=1.1in}
\usepackage{setspace}
\onehalfspacing
\usepackage[authoryear,round]{natbib}
\usepackage{booktabs}
\usepackage{tabularx}
\usepackage{array}
\usepackage{microtype}
\usepackage{xurl}
\usepackage{hyperref}
\hypersetup{colorlinks=false,pdfborder={0 0 0}}
\setlength{\emergencystretch}{3em}
\newcolumntype{Y}{>{\raggedright\arraybackslash}X}
\graphicspath{{../data/processed/scene_networks/figures/}{../data/processed/genre_analysis/}{../data/processed/scene_networks/diffusion/figures/}}

\title{Local scene emergence in heavy metal}
\author{David Chivers\thanks{Department of Economics, Durham University, \href{mailto:david.chivers@durham.ac.uk}{david.chivers@durham.ac.uk}.}}
\date{\today}

\begin{document}

\maketitle

\begin{abstract}
This paper studies local scene emergence in heavy metal. Using the full Metal Archives band census
and member dump, I build a city-level collaboration dataset from 971,877 musician-band edges
covering 682,126 musicians and link it to a dynamic city-by-genre panel of local scene emergence.
A city-genre cell is counted as emergent in the year it first crosses the paper's operational
fifth-band threshold in that broad genre family. In the preferred exact-year fixed-effects sample,
the panel contains 171,653 city-genre-years and 5,306 emergence events. Two facts organize the
paper. First, exact-year emergence rates rise sharply with local band stock and with the depth of
the local multi-band musician pool. Second, in the preferred exact-year transition specification,
later transitions into operational scene status track three stable lagged scene predictors: local
spawning flow, target-genre active bands, and target-genre multi-band musician depth. These
patterns survive nearby threshold changes and geography-clean restrictions, but attenuate sharply
when the risk set is restricted to city-genre cells still far below the threshold. The paper is
therefore strongest on late-stage local niche consolidation rather than on the earliest seed stage
of scene formation. The contribution is reduced-form and measurement-driven rather than causal.
Heavy metal is useful because it makes local project overlap and worker mobility unusually
observable, allowing the paper to study how localized capability, embedded project formation, and
within-niche recombination are associated with local scene emergence.
\end{abstract}

\section{Introduction}
This paper studies local scene emergence in heavy metal. The question is not why some cities have
more bands in general. It is why some city-genre environments move from scattered local entry to a
recognizable local scene earlier than others. I treat that process as a localized production
problem in a project-based creative industry, where bands are projects, musicians move across
teams, and new local niches are easier to sustain once enough specialized capability has
accumulated.

That question is difficult to study with the usual retrospective approach. Famous scenes are often
identified only after genre labels are settled, canonical bands are known, and local histories can
be written in hindsight. But genre recognition is not the same thing as genre formation. By the
time a label is stabilized, much of the relevant local infrastructure may already have formed
through repeated collaboration, spinouts, and worker mobility across projects. The empirical
challenge is therefore to measure local scene formation before the historical narrative fully
settles.

Heavy metal offers an unusually useful setting for that task. Metal Archives provides a broad band
census with genre strings, formation years, and location information, while the member dump makes
it possible to observe overlapping personnel ties across bands. I use those sources to build a
city-level collaboration dataset from 971,877 musician-band edges covering 682,126 musicians, and
I link that object to a dynamic city-by-genre panel of local scene emergence. The resulting panel
contains 56,978 city-genre cells, of which 5,756 eventually cross the paper's operational
emergence threshold. The preferred exact-year fixed-effects sample contains 171,653
city-genre-years and 5,306 emergence events.

The paper's outcome is intentionally operational. A city-genre cell is counted as emergent in the
year it first crosses the paper's fifth-band threshold in that broad genre family. This is not a
claim that genre identity literally begins at band number five. It is a transparent empirical rule
for the transition from isolated local entry to operational scene status. That transparency is a
feature of the design, but it also defines the paper's main empirical risk: one must show that the
results are not merely a mechanical consequence of thicker places crossing a stock threshold
sooner.

Two empirical facts organize the paper. The first is descriptive. Exact-year emergence rates rise
sharply with local band stock and with the depth of the local multi-band musician pool. The second
is reduced-form. In exact-year fixed-effects transition specifications, later city-genre
transitions into operational scene status track three stable lagged predictors: local spawning
flow, target-genre active bands, and target-genre multi-band musician depth. Composition-heavy
extensions and broader connector measures appear only in background work because they do not
improve the main paper object. Threshold-response checks show that this pattern is not pinned to
the fifth-band cutoff, but they also narrow the claim: the paper is strongest on the later
approach to emergence rather than on the earliest seed stage of scene formation.

The contribution is therefore narrower than the broadest version of the project, but clearer.
First, the paper introduces a large city-level collaboration dataset for metal built from the full
band census and member records rather than from a curated set of famous scenes. Second, it treats
genre emergence as a measurable local event that can be studied in panel form before retrospective
scene histories dominate the evidence. Third, it shows in that panel that later transitions into
operational scene status are most closely associated with local spawning, focal-niche band stock,
and focal-niche overlapping-worker depth, with the clearest evidence coming near the later
approach to the scene threshold.

The paper remains reduced-form rather than causal. It does not identify an exogenous source of
scene formation, observe knowledge transfer directly, or date the true historical origin of a
genre. The contribution is instead a carefully validated empirical pattern. In a large panel of
not-yet-emerged city-genre cells, transitions into operational scene status are more likely where
local capability is thicker, project spinouts are more common, and workers are already recombining
within related local projects. Heavy metal is useful here not because it is culturally
exceptional, but because the setting makes project overlap and worker mobility unusually
observable.

The rest of the paper proceeds as follows. Section 2 situates the paper in the related
literature. Section 3 presents the empirical motivation. Section 4 sets out a compact model of
local niche formation. Section 5 describes the data and sample construction. Section 6 presents the
empirical approach. Section 7 reports the descriptive and reduced-form results. The appendix
collects the detailed sample-construction and variable-definition objects behind the main
specification.


\section{Related literature}
This paper speaks most directly to the economic geography of creative industries and to the small
literature on music scenes as localized production systems. \citet{scott_1999} provides the
classic economic-geography statement that recorded music production is concentrated in dense local
environments with specialized workers, organizations, and complementary institutions. The closest
empirical papers are \citet{klement_strambach_eg_2019} and
\citet{klement_strambach_rs_2019}. They show that related variety matters for urban music-scene
creativity and that new genres emerge from diversification processes that remain strongly local
even when extra-local influences matter. My paper fits squarely in that tradition, but shifts the
evidence base away from case studies and retrospective scene histories toward a large
city-by-genre panel built from linked band and member data.

The paper also contributes to a broader economics literature on agglomeration, local externalities,
and the formation of new varieties. \citet{duranton_puga_2004} provide the core micro-foundations:
dense places create advantages through sharing, matching, and learning. \citet{glaeser_etal_1992}
link local variety and dynamic externalities to urban growth, while \citet{storper_venables_2004}
stress the continuing value of face-to-face contact and local buzz even when communication
technologies improve. \citet{moretti_2004} provides direct evidence that local stocks of skilled
workers generate spillovers. These papers are not about music, but they provide the right
economics language for the paper's central claim: local scene thickness can lower the cost of
forming new creative projects because matching is easier, tacit know-how is more available, and
the local capability base is deeper.

The closest mechanism papers are on worker mobility, spinouts, and recombination.
\citet{franco_filson_2006} study spinouts and employee mobility as channels through which
organizational capabilities propagate, \citet{moen_2000} asks whether technical personnel mobility
is a source of R\&D spillovers, and \citet{weitzman_1998} provides the cleanest formal logic for
economically meaningful recombination. Bands are not firms in the standard sense, but the mapping
is useful: workers accumulate project-specific know-how, some leave to found new projects, and
overlapping experience across projects expands the scope for productive recombination. The paper's
spawning-flow and multi-band-musician variables are reduced-form traces of those mechanisms.

At a more abstract level, the paper uses a product-variety interpretation of genre emergence.
\citet{dixit_stiglitz_1977} and \citet{romer_1990} show why the creation of new varieties can be
economically meaningful even when the relevant outcome is not technological improvement in the
narrow sense. The point is not that metal scenes literally replicate aggregate growth theory. It is
that a new niche can be understood as a differentiated product environment that becomes locally
viable once enough project formation, capability, and specialization accumulate.

The paper is also informed by the economics of music markets under digitization.
\citet{ferreira_waldfogel_2013} show that geography and language continue to structure world music
trade, while \citet{waldfogel_2015} shows that digitization changed the supply side of recorded
music by lowering production and distribution frictions and expanding entry. These papers matter
mainly as boundary conditions. The question here is not whether digital technology eliminated
geography in music, but whether some places remain better at generating and sustaining local niche
production even in a more digital environment.

Finally, the paper's outcome is informed by work on genres as social forms rather than fixed
technical taxonomies. \citet{lena_peterson_2008} is especially important here. Genres are social
classifications with trajectories, organizational supports, and shifting boundaries, not objective
technologies with a single uncontested birth date. For that reason, the paper does not claim to
date the true global origin of a genre. Instead, it studies when a local city-genre niche becomes
thick enough to support repeated entry and collaboration. The empirical motivation in the next
section shows those forces in the raw data, and the model that follows compresses them into a
deliberately narrow object: a threshold-crossing entry problem driven by local niche thickness,
embedded spinouts, and recombinant worker depth.


\section{Empirical motivation}
The broader empirical setting is one of rapid expansion and differentiation. Figure
\ref{fig:genre_proliferation} shows the growth of major metal genre families over time. The point
is not that every local scene follows the same path. It is that the paper studies an expanding and
internally differentiating creative field rather than one isolated local anomaly. That makes it
natural to ask why some places form recognizable local scenes earlier than others.

\begin{figure}[htbp]
\centering
\includegraphics[width=0.92\textwidth]{top_genre_families_over_time.png}
\caption{Global expansion and differentiation of major metal genre families.}
\label{fig:genre_proliferation}
\end{figure}

The same pattern is visible inside the genre taxonomy itself. Figure \ref{fig:subgenre_tree}
summarizes the branching structure of metal subgenres in the data. The horizontal axis records the
first year in which each subgenre appears in the Metal Archives-derived band panel, while the
vertical layout is only a visualization device. This should be read as an illustrative label tree,
not a literal genealogy. Its role in the paper is narrower: it shows that later subgenres branch
out of earlier trunks over time, so the empirical setting is one of cumulative variety creation
rather than repeated entry into a fixed category.

\begin{figure}[htbp]
\centering
\includegraphics[width=\textwidth]{subgenre_split_tree_upward.png}
\caption{Branching tree of metal subgenres. The horizontal axis records the first year each subgenre is observed in the Metal Archives-derived band data.}
\label{fig:subgenre_tree}
\end{figure}

The scene question is also spatial. Once a genre becomes legible in a few early hubs, later scenes
appear in other cities and countries. Figure \ref{fig:black_metal_diffusion} illustrates that
pattern for black metal. It does not identify a causal diffusion process and should not be read as
a single-birthplace story. Its role is simply to show that the relevant empirical object is a
sequence of local scene formations rather than one isolated national event.

\begin{figure}[htbp]
\centering
\includegraphics[width=\textwidth]{black_metal_diffusion_snapshots.png}
\caption{Black metal diffusion snapshots in 1990 and 2000. Marker area scales with the number of active local black-metal bands in each city-year, with square-root scaling for readability.}
\label{fig:black_metal_diffusion}
\end{figure}

Taken together, these figures do only two pieces of setup work. First, heavy metal generates enough
entry and differentiation to study local niche formation as an economic outcome. Second, that
process is uneven across places and unfolds through repeated local scene formation. The empirical
analysis therefore turns to the city-by-genre panel and asks which local scene characteristics
predict the transition into operational scene status. The model section takes that narrower
threshold-crossing problem, rather than the global origin of a genre, as the object to explain.

\clearpage

\section{Model}
The literature and descriptive evidence point to a narrower theoretical task than a full model of
genre origin or cultural diffusion. The paper needs a framework that explains why some
city-genre cells are more likely than others to cross from thin local activity into repeated local
entry. In the data, that transition is most strongly associated with local spawning flow,
target-genre active bands, and target-genre multi-band musician depth.

The empirical object in this paper is therefore local niche formation in a project-based creative
industry. Bands are projects, musicians are skilled workers, and a local scene emerges when
activity in a city-genre cell becomes thick enough to sustain repeated entry. The model is not a
theory of technological progress in the usual sense. It is a compact model of local variety
creation in which entry depends on local niche thickness, spinouts by already embedded
participants, and recombination through overlapping workers.

\subsection*{Environment}

Time is discrete. Cities are indexed by $c$, genre niches by $g$, and years by $t$. For a
city-genre cell that is still at risk of emergence, the local state at the start of period $t$ is
summarized by
\[
\left(B_{cg,t-1},\; M_{cg,t-1},\; \rho_{c,t-1}\right),
\]
where $B_{cg,t-1}$ is the number of active local bands in the focal genre, $M_{cg,t-1}$ is the
number of local multi-band musicians active in that genre, and $\rho_{c,t-1}$ is the share of
potential founders who are already embedded in the broader local scene. A band is one
differentiated local variety, and a scene emerges when the local niche becomes thick enough to
sustain repeated entry rather than isolated projects.

\subsection*{Entry}

In each period, a mass of potential founders considers launching a new band in genre $g$ in city
$c$. Founder $i$ draws embedded status $e_i \in \{0,1\}$ with
$\Pr(e_i = 1) = \rho_{c,t-1}$, an idiosyncratic fixed cost $f_i$, and a payoff shock
$\varepsilon_i$. Expected net entry value is
\begin{equation}
\Pi_{icgt}
=
\bar{\pi}_g
+ \phi_B \log(1 + B_{cg,t-1})
+ \phi_M \log(1 + M_{cg,t-1})
+ \gamma e_i
- f_i
+ \varepsilon_i,
\qquad
\phi_B,\phi_M,\gamma > 0.
\end{equation}

The term in $B_{cg,t-1}$ captures local niche thickness and labor pooling. The term in
$M_{cg,t-1}$ captures recombinant capacity through overlapping workers. The embedded-founder
premium $\gamma$ captures inherited local know-how from prior participation in the scene. Broader
city conditions are left in the background here and are absorbed empirically by fixed effects and
the paper's broader scene controls rather than being modeled explicitly.

Entry occurs when $\Pi_{icgt} \ge 0$. Aggregate entry is
\begin{equation}
N_{cgt}
=
\int \mathbf{1}\{\Pi_{icgt} \ge 0\} \, dH_{ct}(i),
\end{equation}
where $H_{ct}$ is the distribution of founder types in city $c$ at date $t$. It follows
immediately that
\[
\frac{\partial N_{cgt}}{\partial B_{cg,t-1}} > 0,
\qquad
\frac{\partial N_{cgt}}{\partial M_{cg,t-1}} > 0,
\qquad
\frac{\partial N_{cgt}}{\partial \rho_{c,t-1}} > 0.
\]

\subsection*{Dynamics and emergence}

Band stock evolves as
\begin{equation}
B_{cgt} = (1 - \delta) B_{cg,t-1} + N_{cgt},
\qquad
\delta \in (0,1).
\end{equation}

Let $N^E_{cg't}$ denote entry by embedded founders in genre $g'$. The observed spawning variable is
\begin{equation}
S_{ct} = \sum_{g'} N^E_{cg't}.
\end{equation}
Because $S_{ct}$ is indexed at the city-year level rather than the city-genre level, it is best
read as a proxy for broad local spinout intensity rather than as a purely focal-genre count.

A local scene emerges when activity crosses a viability threshold:
\begin{equation}
E_{cgt} = \mathbf{1}\{B_{cgt} \ge \bar{B}\},
\qquad
\bar{B} = 5.
\end{equation}
This threshold is operational rather than ontological. It marks the point at which a local niche
looks thick enough to sustain repeated entry, not the literal historical birth of a genre.

For cells that are still at risk, the one-period emergence hazard is
\begin{equation}
h_{cgt} = \Pr(E_{cgt} = 1 \mid E_{cg,t-1} = 0).
\end{equation}
Because $B_{cgt}$ depends on prior niche stock and current entry, this hazard is increasing in
$B_{cg,t-1}$, $M_{cg,t-1}$, and the latent embedded-founder share $\rho_{c,t-1}$.

\subsection*{Empirical mapping}

The model maps directly to the paper's three headline predictors:
\begin{align*}
B_{cg,t-1} &\leftrightarrow \text{target-genre active bands}, \\
M_{cg,t-1} &\leftrightarrow \text{target-genre multi-band musicians}, \\
\rho_{c,t-1} &\leftrightarrow \text{embedded-founder margin, proxied by } S_{c,t-1}.
\end{align*}

The reduced-form exact-year specification is therefore a compact empirical analogue of the
threshold-crossing condition:
\begin{equation}
\Pr(E_{cgt}=1 \mid E_{cg,t-1}=0)
=
G\!\left(
\alpha_{cg}
+ \lambda_t
+ \beta_S S_{c,t-1}
+ \beta_B B_{cg,t-1}
+ \beta_M M_{cg,t-1}
\right),
\end{equation}
where $G(\cdot)$ is increasing. The estimated specification adds broader scene controls and
smoothed lags, but this compact version captures the paper's core economic object.

\subsection*{Predictions and scope}

The model yields three predictions aligned with the main results. First, thicker target-genre
stock should raise later emergence because deeper local niche activity lowers the effective cost of
entry. Second, deeper overlapping-worker pools should raise later emergence because recombination
and team assembly are easier where more musicians already work across related projects. Third,
higher city-year spawning flow should predict later emergence because places that generate more
embedded spinouts also generate more viable founders than de novo entry alone would.

The interpretation remains deliberately modest. The paper does not observe knowledge transfer
directly, does not treat the emergence threshold as structural, and does not claim that spawning
is a purely genre-specific causal mechanism. The model instead rationalizes a reduced-form pattern:
new local scenes are more likely to emerge where local labor pooling is thicker, embedded
participants generate more new projects, and overlapping workers are already recombining across
related bands.


\section{Data}
The empirical analysis combines a complete band reference, a complete member file, and a set of
derived city-level collaboration panels. The cleaned all-metal reference built from the full Metal
Archives snapshot contains 163,965 bands, of which 131,511 carry a usable formed year. The member
dump contributes 1,038,533 raw member-role rows, which collapse to 971,877 unique
musician-band edges once duplicate within-band role rows are removed. These source files make it
possible to move beyond anecdotal scene histories and study local collaboration at scale.

The data construction has two substantive steps. The first builds yearly local collaboration
snapshots by linking bands that share at least one musician and are active in the same year. These
city-year graphs are then summarized into measures of local scale, organizational variety, and
cross-project worker mobility. The second builds the city-by-genre emergence panel. Genre strings
are mapped into broad genre families, bands are assigned to cities using the cleaned location
field, and each city-genre cell is dated by its first observed band and its emergence year. A cell
is counted as emergent when the city reaches its fifth band in that broad genre family.

The preferred estimation sample is the roll-5 exact-year panel. It contains 171,653 city-genre-
year observations and 5,306 emergence events. Table \ref{tab:sample_construction} summarizes the
construction of that sample from the raw band and member universes to the final estimating panel.
Appendix Table \ref{tab:appendix_sample_construction} reports the fuller construction audit,
including the intermediate network objects, the geography exclusions, and the nearby roll-3
robustness sample. Appendix Tables \ref{tab:appendix_measurement_validation} and
\ref{tab:appendix_historical_sanity} report the measurement-validation objects behind the live
emergence definition.

\begin{table}[htbp]
\centering
\caption{Sample construction for the preferred city-genre scene panel}
\label{tab:sample_construction}
\small
\begin{tabularx}{\textwidth}{@{}lYcY@{}}
\toprule
Stage & Object & Count & Restriction or construction rule \\
\midrule
1 & Raw all-metal Metallum band universe & 163,965 & Upstream cleaned all-metal reference \\
2 & Bands with standardized countries & 161,945 & Country standardization succeeds \\
3 & Bands with usable formed year & 131,511 & Parsable formed year required \\
4 & Raw member-role universe & 1,038,533 & Delivered member dump before collapse \\
5 & Unique musician-band edges & 971,877 & Duplicate role rows collapsed to one band-musician edge \\
6 & Unique musicians & 682,126 & Unique musician keys after edge collapse \\
7 & Reference-matched bands in edge file & 156,729 & Band ID matched back to cleaned band reference \\
8 & City-genre cells in emergence panel & 56,978 & Usable city, usable formed year, mapped genre family \\
9 & City-genre cells reaching emergence & 5,756 & City reaches its fifth band in the genre family \\
10 & Preferred exact-year FE sample & 171,653 & Valid lagged roll-5 predictors and at-risk city-genre-years \\
11 & Preferred emergence events & 5,306 & Snapshot year equals emergence year within preferred sample \\
\bottomrule
\end{tabularx}

\vspace{0.4em}
\parbox{0.96\textwidth}{\footnotesize Notes: A city-genre cell is counted as emergent in the year
it reaches its fifth band in the broad genre family. The preferred estimating sample is the roll-5
exact-year fixed-effects panel used in the current main specification.}
\end{table}



\section{Empirical approach}
The unit of observation is the city-by-genre-by-year cell. The outcome is an indicator for whether
that cell emerges in a given year. The headline regressors are local spawning flow, target-genre
active bands, and target-genre multi-band musicians. The preferred specification also includes
broader city-level controls for overall active bands, nontrivial communities, and total multi-band
musicians. All predictors are transformed as one-year-lagged rolling averages and z-scored within
the estimation sample. Table \ref{tab:main_variables} reports the main variables used in the
current draft. Appendix Table \ref{tab:appendix_variable_crosswalk} reports the fuller variable
crosswalk, including the code-side variable names, units of observation, and the level at which
the rolling transformations are constructed before the full panel merge. This standardization is
useful later in the results section because the coefficients can be compared on a common
within-sample scale: each coefficient is the change in emergence probability associated with a
one-standard-deviation increase in the lagged predictor.

\begin{table}[htbp]
\centering
\caption{Main variables in the exact-year scene specification}
\label{tab:main_variables}
\small
\begin{tabularx}{\textwidth}{@{}lY Y@{}}
\toprule
Variable & Definition & Transformation and role \\
\midrule
Emergence this year & Indicator equal to one when the snapshot year is the emergence year for a city-genre cell & Binary outcome in exact-year fixed effects \\
Local spawning flow & Bands formed in a city-year whose founding cohort includes at least one musician with prior local band history in that city & Log transformed, one-year-lagged roll-5 average, z-scored; headline mechanism \\
Target-genre active bands & Active local bands in the focal broad genre family & Log transformed, one-year-lagged roll-5 average, z-scored; headline mechanism \\
Target-genre multi-band musicians & Local musicians active in the focal genre who hold at least two active local bands in that same genre & Log transformed, one-year-lagged roll-5 average, z-scored; headline mechanism \\
Overall active bands & Active local bands of any genre in the city-year snapshot & Log transformed, one-year-lagged roll-5 average, z-scored; control \\
Nontrivial communities & Number of detected communities with at least two active bands in the local yearly network & One-year-lagged roll-5 average, z-scored; control \\
Overall multi-band musicians & Active local musicians with at least two active local bands of any genre & Log transformed, one-year-lagged roll-5 average, z-scored; control \\
\bottomrule
\end{tabularx}

\vspace{0.4em}
\parbox{0.96\textwidth}{\footnotesize Notes: The preferred specification uses city-genre and year
fixed effects with standard errors clustered by city. All reported continuous predictors are
z-scored within the estimation sample.}
\end{table}


The preferred exact-year specification is a reduced-form transition equation for the probability
that a city-genre cell that has not yet emerged first crosses the paper's operational scene
threshold in year $t$:
\[
\Pr(E_{cgt}=1 \mid E_{cg,t-1}=0) = F\left(\alpha_{cg} + \lambda_t + \beta_1 S_{c,t-1}
 + \beta_2 B_{cg,t-1} + \beta_3 M_{cg,t-1}\right),
\]
where $E_{cgt}$ indicates whether city $c$ and genre $g$ emerge in year $t$, $S_{c,t-1}$ is
lagged local spawning flow, $B_{cg,t-1}$ is lagged target-genre active band stock, and
$M_{cg,t-1}$ is lagged target-genre multi-band musician depth. The design absorbs city-genre fixed
effects and year fixed effects, with standard errors clustered by city. This specification is
useful because it compares the same city-genre environment over time rather than relying on static
differences between famous and obscure scenes. It should therefore be read as a within-cell
transition equation on the risk set, not as a structural production function for scenes.

The design should not be read as delivering causal identification. Stronger local scenes may
generate some of the measured predictors, and omitted local shocks may move both current scene
conditions and later emergence. The coefficients therefore support a disciplined statement about
which local scene features predict later transitions into operational scene status, not a claim
that the paper directly observes knowledge transfer or the exact historical birth date of a genre.


\section{Results}
The first result is descriptive. Exact-year emergence rates rise sharply with local scene
thickness. When a city has fewer than 10 active local bands, the exact-year emergence rate is
1.4 percent. When it has 40 or more, the rate rises to 4.0 percent. The same pattern appears for
multi-band depth: emergence rises from 1.4 percent with 0--1 multi-band musicians to 3.7 percent
with 10 or more. Figure \ref{fig:critical_mass} reports this critical-mass pattern directly.

\begin{figure}[htbp]
\centering
\includegraphics[width=0.92\textwidth]{scene_critical_mass_event_rates.png}
\caption{Exact-year emergence rates by local scene thickness and multi-band depth.}
\label{fig:critical_mass}
\end{figure}

The main regression result appears in the preferred exact-year fixed-effects transition
specification. The counts-only roll-5 design is the cleanest current main-table object because it
isolates the stable unconditional scene terms. The dependent variable is whether a city-genre cell
that has not yet emerged crosses the operational fifth-band threshold in the current year.
Table \ref{tab:preferred_results} reports the preferred paired presentation: a core counts-only
panel and a composition extension. The main paper read comes from Panel A. Because the lagged
predictors are z-scored, the coefficients are directly comparable on a common within-sample scale.
Local spawning flow enters at 0.0134, target-genre active bands at 0.0505, and target-genre
multi-band musicians at 0.0227. Relative to the sample event rate of 3.09 percent, those are
economically large associations on the linear-probability scale: a one-standard-deviation increase
in local spawning flow corresponds to about 1.3 percentage points higher transition probability, a
one-standard-deviation increase in target-genre active bands corresponds to about 5.1 percentage
points, and a one-standard-deviation increase in target-genre multi-band depth corresponds to
about 2.3 percentage points. The largest coefficient is therefore not generic city size, but
focal-niche thickness.

\begin{table}[htbp]
\centering
\caption{Preferred exact-year fixed-effects scene results}
\label{tab:preferred_results}
\small
\begin{tabular}{@{}lcc@{}}
\toprule
Predictor & Panel A: core counts-only & Panel B: composition extension \\
\midrule
Overall active bands & 0.0135*** & 0.0174*** \\
 & (0.0043) & (0.0044) \\
Nontrivial communities & 0.0062*** & 0.0008 \\
 & (0.0018) & (0.0020) \\
Overall multi-band musicians & -0.0223*** & -0.0244*** \\
 & (0.0028) & (0.0030) \\
Local spawning flow & 0.0134*** & 0.0224*** \\
 & (0.0017) & (0.0025) \\
Spawn share &  & -0.0066*** \\
 &  & (0.0012) \\
Target-genre active bands & 0.0505*** & 0.0488*** \\
 & (0.0015) & (0.0017) \\
Target-genre broker musicians &  & 0.0579*** \\
 &  & (0.0028) \\
Target-genre broker share &  & -0.0724*** \\
 &  & (0.0022) \\
Target-genre multi-band musicians & 0.0227*** & 0.0331*** \\
 & (0.0015) & (0.0027) \\
Target-genre multi-band share &  & -0.0195*** \\
 &  & (0.0023) \\
\midrule
$N$ & 171,653 & 171,653 \\
Events & 5,306 & 5,306 \\
Event rate & 0.0309 & 0.0309 \\
$R^2$ & 0.0422 & 0.0581 \\
Fixed effects & City-genre + year & City-genre + year \\
Clustering & City & City \\
\bottomrule
\end{tabular}

\vspace{0.4em}
\parbox{0.96\textwidth}{\footnotesize Notes: Standard errors in parentheses. All predictors are
one-year-lagged \texttt{roll-5} averages z-scored within the estimation sample, so coefficients are
comparable as within-sample percentage-point changes in emergence probability from a
one-standard-deviation increase in the predictor. *** $p<0.01$. Panel A is the active headline
specification. Panel B is retained as a background composition extension because the share terms
are substantively informative even though the main paper read comes from the core counts-only
panel. The negative coefficient on overall multi-band musicians should be read as a conditional
control result: holding focal-genre depth fixed, broader city-wide overlapping labor appears less
helpful than overlap concentrated within the target niche.}
\end{table}


The control pattern is also interpretable once the table is read conditionally rather than
literally. Overall active bands remain positive, but overall multi-band musicians enter negatively.
That sign is not a claim that worker overlap is bad for scenes. It is the opposite comparison.
Holding fixed focal-genre active bands and focal-genre multi-band depth, broader city-wide
multi-band overlap appears to proxy for labor spread across many niches rather than for depth in
the focal one. The paper's mechanism read is therefore about niche-specific recombination, not
generic scene-wide overlap. This is also why the positive target-genre multi-band coefficient is
the more substantively important labor-pool margin.

The direct threshold-response package is reported in Appendix Tables
\ref{tab:appendix_threshold_response} and \ref{tab:appendix_threshold_specificity}. The first
result is encouraging. The coefficient pattern is not pinned to the fifth-band cutoff: across
alternative emergence thresholds of three, four, six, and eight bands, local spawning flow,
target-genre active bands, and target-genre multi-band musicians all remain positive. The harder
risk-set exercises are more sobering. When I drop observations one band below the fifth-band
cutoff, all three headline terms remain positive but become much smaller. When I keep only
city-genre cells that are still well below the cutoff, the coefficients attenuate sharply and the
multi-band term is no longer distinguishable from zero. The disciplined interpretation is
therefore narrower than the most expansive version of the project. The paper is strongest on
late-stage local niche consolidation rather than on the earliest seed stage of scene formation.
A focal labor-pool specificity check still supports the recombination margin: conditional on
target-genre active bands and total target-genre active musicians, target-genre multi-band
musicians remain positive while total target-genre active musicians turn negative.

The measurement-validation package is reported in Appendix Tables
\ref{tab:appendix_measurement_validation} and \ref{tab:appendix_historical_sanity}. The current
preferred sample is already geography-clean on the project's audit rule: region-like and malformed
city labels are excluded upstream, so none survive into the live fixed-effects panel. Tightening
the sample further does not materially change the main coefficients. Dropping the small dense
labels flagged for closer inspection, restricting the sample to `1990` onward, and restricting to
the high-coverage countries all leave the three headline terms close to their baseline values.
The historical sanity cases are also reassuring. Operational emergence dates such as Birmingham
heavy metal in `1977`, Tampa death metal in `1987`, Bergen black metal in `1992`, and Gothenburg
melodic death metal in `1996` are at least broadly plausible. The Bay Area case is useful for a
different reason: its raw timing looks plausible too, but the label is explicitly excluded because
it names a wider unit rather than a city.

The coefficient pattern points to a specific view of local creative development. City-genre cells
are more likely to transition into operational scene status when three conditions hold. First,
cities generate new projects through local spinouts. Second, the focal genre already has a thicker
local stock of active bands. Third, musicians in that focal genre work across multiple local
projects, which is consistent with labor pooling and recombination within the niche.
Figure \ref{fig:coefficients} visualizes the current headline coefficients.

\begin{figure}[htbp]
\centering
\includegraphics[width=0.78\textwidth]{scene_preferred_mechanism_coefficients.png}
\caption{Preferred core exact-year fixed-effects coefficients on z-scored lagged predictors.}
\label{fig:coefficients}
\end{figure}

Background role-composition checks add little to this interpretation. Only guitarist share among
switchers shows any visible signal, and even that is weakly negative rather than positive
($-0.0013$, $p = 0.062$); drums, bass, vocals, and keyboards are all close to zero. I therefore
leave role composition out of the headline narrative. The stable margin is broader
overlapping-worker depth within the focal niche, not one privileged instrument class.

The remaining digital-era checks are secondary and do not change the paper's main read. A broad
period split suggests that local spawning flow and target-genre active-band stock attenuate after
`2005`, while the target-genre multi-band-musician coefficient remains stable. That is useful as a
descriptive nuance, but it is not central to the paper's contribution. I therefore leave the
broader internet-adoption and city-broadband probes to Appendix Section \ref{app:digital_checks}.
Those exercises are either too broad, too late, or too small to identify a local digital
mechanism cleanly, and none of them overturn the main scene result.

The contribution is reduced-form rather than causal. That is still enough to establish a clear
empirical pattern. Thicker local capability, stronger spinout flow, and deeper within-genre
overlapping experience are all associated with later scene emergence. Heavy metal is useful here
not because it is culturally exceptional, but because the setting makes local project overlap and
worker mobility unusually observable. In that sense, the paper speaks to a broader economic
question: how localized capability and recombination shape the emergence of new creative
varieties.


\appendix
\setcounter{table}{0}
\renewcommand{\thetable}{A\arabic{table}}
\renewcommand{\theHtable}{appendix.\thetable}
\setcounter{figure}{0}
\renewcommand{\thefigure}{A\arabic{figure}}
\renewcommand{\theHfigure}{appendix.\thefigure}
\clearpage
\section{Appendix material for the main scene specification}

This appendix gathers the audit objects that sit behind the active main-text scene result. The
goal is not to reopen the design, but to make the preferred sample and variable definitions
transparent enough that the empirical package is fully auditable.

\subsection{Sample construction}

Appendix Table \ref{tab:appendix_sample_construction} reports the full construction of the
city-genre panel from the raw band and member universes to the preferred \texttt{roll-5}
exact-year fixed-effects sample. It also records the geography-audit exclusions and the nearby
\texttt{roll-3} robustness sample so that the main-text estimation object is easy to locate inside
the broader data pipeline.

\begin{table}[p]
\centering
\caption{Full sample construction for the city-genre scene panel}
\label{tab:appendix_sample_construction}
\tiny
\setlength{\tabcolsep}{3pt}
\renewcommand{\arraystretch}{1.08}
\begin{tabularx}{\textwidth}{@{}>{\raggedright\arraybackslash}p{0.06\textwidth}Y>{\raggedleft\arraybackslash}p{0.12\textwidth}YY@{}}
\toprule
Stage & Object & Count & Restriction or construction rule & Source \\
\midrule
1 & Raw all-metal Metallum band universe & 163,965 & Upstream cleaned all-metal reference. & All-metal summary \\
2 & Bands matched to standardized countries & 161,945 & Country standardization succeeds. & All-metal summary \\
3 & Bands with usable formed year & 131,511 & Parsable formed year required. & All-metal summary \\
4 & Raw member-role universe & 1,038,533 & Delivered member dump before collapse. & Member ingest summary \\
5 & Unique musician-band edges & 971,877 & Duplicate role rows collapsed to one band-musician edge. & Member ingest summary \\
6 & Unique musicians & 682,126 & Unique musician keys after edge collapse. & Member ingest summary \\
7 & Reference-matched bands in edge file & 156,729 & Band ID matched back to cleaned band reference. & Member ingest summary \\
8 & Cities with any member-linked bands & 22,572 & Requires usable city label in band metadata. & City network summary \\
9 & Eligible cities for descriptive network baseline & 2,319 & At least 10 bands in city. & City network summary \\
10 & City-genre cells in emergence panel & 56,978 & Usable city, usable formed year, mapped genre family. & City-genre emergence summary \\
11 & City-genre cells reaching emergence threshold & 5,756 & City reaches its fifth band in the genre family. & City-genre emergence summary \\
12 & Live city-year network snapshots & 32,684 & Strict timing baseline with minimum active city bands of 10. & Scene cleaning audit \\
13 & Live city-genre timing rows & 5,429 & Current strict city baseline after geography audit. & Scene cleaning audit \\
14 & Preferred exact-year FE sample (roll-5) & 171,653 & Non-missing first-band year, snapshot year on or after first-band year, pre-emergence or emergence-year risk set, and lagged roll-5 predictors available. & Preferred mechanism table \\
15 & Preferred exact-year FE emergence events (roll-5) & 5,306 & Snapshot year equals emergence year within preferred sample. & Preferred mechanism table \\
\bottomrule
\end{tabularx}

\vspace{0.4em}
\parbox{0.96\textwidth}{\footnotesize Notes: The preferred paper sample is the roll-5 exact-year
fixed-effects panel. The nearby roll-3 robustness panel contains 180,498
city-genre-years and 5,422 emergence events. The geography audit is a reusable exclusion register
rather than a nested sample row: its audit universe contains 2,443 labels, with 86
region-like labels and 66 missing-or-multi-place labels, for 152
city-baseline exclusions. A city-genre cell is counted as emergent in the year it reaches its
fifth band in the broad genre family.}
\end{table}


\subsection{Variable definitions}

Appendix Table \ref{tab:appendix_variable_crosswalk} defines the variables used in the active
headline specification. It distinguishes city-year controls from city-genre mechanisms, records
where the rolling transformations are applied, and keeps the paper-facing labels tied to the code
variables used in the estimation pipeline.

\begin{table}[p]
\centering
\caption{Variable definitions for the main scene specification}
\label{tab:appendix_variable_crosswalk}
\tiny
\setlength{\tabcolsep}{3pt}
\renewcommand{\arraystretch}{1.08}
\begin{tabularx}{\textwidth}{@{}>{\raggedright\arraybackslash}p{0.14\textwidth}>{\raggedright\arraybackslash}p{0.20\textwidth}>{\raggedright\arraybackslash}p{0.12\textwidth}YY@{}}
\toprule
Paper label & Code variable & Unit of observation & Construction & Estimating transformation and role \\
\midrule
Emergence this year & emerge\_this\_year & city x genre x year & Indicator equal to 1 when snapshot year equals emergence year; emergence year is the year the city reaches its fifth band in the genre family. & Binary outcome in exact-year fixed effects; outcome variable. \\
Overall active bands & n\_active\_bands \newline $\rightarrow$ log\_active\_bands & city x year & Active local bands in the yearly shared-member network snapshot. & log(1 + x), one-year-lagged roll-5 average, z-scored within the estimation sample; control. \\
Nontrivial communities & n\_nontrivial\_ \newline communities & city x year & Number of detected communities with at least two active bands. & One-year-lagged roll-5 average, z-scored; control. \\
Overall multi-band musicians & n\_multi\_band\_ \newline musicians\_total \newline $\rightarrow$ log\_multi\_band\_ \newline musicians\_total & city x year & Active local musicians with two or more active local bands of any genre. & log(1 + x), one-year-lagged roll-5 average, z-scored; control. \\
Local spawning flow & spawn\_bands\_ \newline formed \newline $\rightarrow$ log\_spawn\_bands\_ \newline formed & city x year & Bands formed in the city-year whose founding cohort includes at least one musician with prior local-band history in that city. & log(1 + x), one-year-lagged roll-5 average, z-scored; headline mechanism. \\
Target-genre active bands & genre\_active\_ \newline bands \newline $\rightarrow$ log\_genre\_active\_ \newline bands & city x genre x year & Active local bands in the focal broad genre family. & log(1 + x), one-year-lagged roll-5 average at the city-genre level, z-scored; headline mechanism. \\
Target-genre multi-band musicians & genre\_multi\_band\_ \newline musicians \newline $\rightarrow$ log\_genre\_multi\_band\_ \newline musicians & city x genre x year & Local musicians active in the focal genre who hold two or more active local bands in that same genre. & log(1 + x), one-year-lagged roll-5 average at the city-genre level, z-scored; headline mechanism. \\
\bottomrule
\end{tabularx}

\vspace{0.4em}
\parbox{0.96\textwidth}{\footnotesize Notes: The preferred specification in the active paper draft
is the core roll-5 exact-year fixed-effects table with city-genre and year fixed effects
and standard errors clustered by city. Z-scores are computed within the estimation sample rather
than in a universal preprocessing file. Overall scene predictors are rolled at the city-year level
before the full panel merge. Target-genre predictors are rolled at the city-genre level before the
merge. Connector or broker variables are not part of the active main-text paper package.}
\end{table}


\subsection{Threshold-response package}

Appendix Table \ref{tab:appendix_threshold_response} is the direct answer to the mechanical-threshold
objection. It first reruns the preferred exact-year fixed-effects specification under nearby
emergence cutoffs of three, four, six, and eight bands. It then holds the fifth-band definition
fixed and asks how much of the predictive pattern remains once the risk set is pushed away from the
immediate approach to the cutoff.

\begin{table}[p]
\centering
\caption{Threshold-response package for the preferred exact-year scene specification}
\label{tab:appendix_threshold_response}
\tiny
\setlength{\tabcolsep}{3pt}
\renewcommand{\arraystretch}{1.08}
\begin{tabularx}{\textwidth}{@{}>{\raggedright\arraybackslash}p{0.29\textwidth}YYY>{\raggedleft\arraybackslash}p{0.10\textwidth}>{\raggedleft\arraybackslash}p{0.08\textwidth}@{}}
\toprule
Sample or threshold & Local spawning flow & Target-genre active bands & Target-genre multi-band musicians & $N$ & Events \\
\midrule
\multicolumn{6}{@{}l}{\textit{Panel A. Alternative emergence thresholds}} \\
Threshold = 3 bands & 0.0152*** & 0.0583*** & 0.0123*** & 133,117 & 6,565 \\
 & (0.0023) & (0.0016) & (0.0020) &  &  \\
Threshold = 4 bands & 0.0132*** & 0.0542*** & 0.0169*** & 154,856 & 5,676 \\
 & (0.0018) & (0.0015) & (0.0017) &  &  \\
Threshold = 5 bands & 0.0134*** & 0.0505*** & 0.0227*** & 171,653 & 5,306 \\
 & (0.0017) & (0.0015) & (0.0015) &  &  \\
Threshold = 6 bands & 0.0127*** & 0.0363*** & 0.0160*** & 180,091 & 3,230 \\
 & (0.0014) & (0.0015) & (0.0011) &  &  \\
Threshold = 8 bands & 0.0103*** & 0.0270*** & 0.0161*** & 190,638 & 2,273 \\
 & (0.0011) & (0.0013) & (0.0010) &  &  \\
\midrule
\multicolumn{6}{@{}l}{\textit{Panel B. Restricted fifth-band risk sets}} \\
Full fifth-band risk set & 0.0134*** & 0.0505*** & 0.0227*** & 171,653 & 5,306 \\
 & (0.0017) & (0.0015) & (0.0015) &  &  \\
Drop rows one band below the cutoff & 0.0046*** & 0.0113*** & 0.0024*** & 154,856 & 1,010 \\
 & (0.0008) & (0.0007) & (0.0006) &  &  \\
Keep only cells with lag cumulative formed $\leq 2$ & 0.0010*** & 0.0019*** & -0.0001 & 133,117 & 156 \\
 & (0.0003) & (0.0002) & (0.0002) &  &  \\
\bottomrule
\end{tabularx}

\vspace{0.4em}
\parbox{0.96\textwidth}{\footnotesize Notes: All specifications use the preferred
\texttt{roll-5} exact-year fixed-effects design with city-genre and year fixed effects and
standard errors clustered by city. Panel A varies only the operational emergence cutoff used to
date the city-genre outcome. Panel B keeps the fifth-band definition fixed and changes the risk set
using lagged cumulative formed bands in the focal city-genre cell. The far-below-threshold row
shows that the strongest predictive content comes from the later approach to emergence rather than
from the earliest seed stage. *** $p<0.01$.}
\end{table}


\subsection{Focal labor-pool specificity}

Appendix Table \ref{tab:appendix_threshold_specificity} adds target-genre active musicians to the
preferred fifth-band specification. This is a falsification-style check on whether the positive
multi-band coefficient is simply standing in for broad focal labor-pool size rather than
overlapping worker depth within the niche.

\begin{table}[p]
\centering
\caption{Focal labor-pool specificity check}
\label{tab:appendix_threshold_specificity}
\small
\setlength{\tabcolsep}{4pt}
\renewcommand{\arraystretch}{1.08}
\begin{tabular}{@{}lc@{}}
\toprule
Predictor & Coefficient \\
\midrule
Target-genre active bands & 0.1533*** \\
 & (0.0035) \\
Target-genre active musicians & -0.1049*** \\
 & (0.0027) \\
Target-genre multi-band musicians & 0.0121*** \\
 & (0.0015) \\
\midrule
$N$ & 171,653 \\
Events & 5,306 \\
$R^2$ & 0.0592 \\
\bottomrule
\end{tabular}

\vspace{0.4em}
\parbox{0.88\textwidth}{\footnotesize Notes: This specification adds lagged target-genre active
musicians to the preferred fifth-band exact-year fixed-effects regression. The point is to ask
whether the overlap term is just a proxy for broad focal labor-pool size. City-genre and year
fixed effects are included and standard errors are clustered by city. The multi-band term remains
positive even after conditioning on total focal active musicians. *** $p<0.01$.}
\end{table}


\subsection{Measurement validation}

Appendix Table \ref{tab:appendix_measurement_validation} validates the emergence object on three
margins. It first reports the current geography audit that defines the city baseline. It then reruns
the preferred exact-year specification after dropping the small-dense labels flagged for closer
inspection, after restricting the sample to `1990` onward, and after restricting to the
high-coverage countries that contribute at least `50` emergence events.

\begin{table}[p]
\centering
\caption{Measurement validation for the city-genre emergence object}
\label{tab:appendix_measurement_validation}
\tiny
\setlength{\tabcolsep}{3pt}
\renewcommand{\arraystretch}{1.08}
\begin{tabularx}{\textwidth}{@{}>{\raggedright\arraybackslash}p{0.30\textwidth}YYY>{\raggedleft\arraybackslash}p{0.10\textwidth}>{\raggedleft\arraybackslash}p{0.08\textwidth}@{}}
\toprule
Object or sample & Local spawning flow & Target-genre active bands & Target-genre multi-band musicians & $N$ & Events \\
\midrule
\multicolumn{6}{@{}l}{\textit{Panel A. Geography audit counts}} \\
Audit-universe city labels &  &  &  & 2,443 &  \\
Region-like labels &  &  &  & 86 &  \\
Missing-or-multi-place labels &  &  &  & 66 &  \\
Current city-baseline exclusions &  &  &  & 152 &  \\
Small dense labels flagged for stricter validation &  &  &  & 177 &  \\
Preferred FE sample rows and events &  &  &  & 171,653 & 5,306 \\
Preferred-sample rows in small dense cities &  &  &  & 4,874 & 166 \\
\midrule
\multicolumn{6}{@{}l}{\textit{Panel B. Geography and coverage robustness}} \\
Preferred baseline & 0.0134*** & 0.0505*** & 0.0227*** & 171,653 & 5,306 \\
 & (0.0017) & (0.0015) & (0.0015) &  &  \\
Drop small dense labels & 0.0133*** & 0.0510*** & 0.0215*** & 166,779 & 5,140 \\
 & (0.0017) & (0.0015) & (0.0015) &  &  \\
Restrict to 1990 onward & 0.0135*** & 0.0481*** & 0.0235*** & 166,685 & 5,071 \\
 & (0.0017) & (0.0014) & (0.0016) &  &  \\
Restrict to countries with at least 50 events & 0.0143*** & 0.0514*** & 0.0213*** & 142,372 & 4,434 \\
 & (0.0019) & (0.0017) & (0.0017) &  &  \\
\bottomrule
\end{tabularx}

\vspace{0.4em}
\parbox{0.96\textwidth}{\footnotesize Notes: Panel A reports the current geography screen behind
the city baseline. Region-like and malformed labels are excluded upstream before the live feature
files are built, which is why the preferred FE sample contains zero such rows. Panel B reruns the
preferred \texttt{roll-5} exact-year fixed-effects specification on narrower geography and
coverage subsets. Standard errors are clustered by city. *** $p<0.01$.}
\end{table}


\subsection{Historical sanity cases}

Appendix Table \ref{tab:appendix_historical_sanity} reports a small set of canonical scenes for
which the operational emergence dates should be broadly plausible. It also includes the Bay Area
thrash-metal row to make the geography logic transparent: the raw timing file dates it, but the
city baseline excludes it because `Bay Area` is a wider unit rather than a city-like label.

\begin{table}[p]
\centering
\caption{Selected historical sanity cases}
\label{tab:appendix_historical_sanity}
\small
\setlength{\tabcolsep}{4pt}
\renewcommand{\arraystretch}{1.08}
\begin{tabularx}{\textwidth}{@{}>{\raggedright\arraybackslash}p{0.24\textwidth}>{\raggedright\arraybackslash}p{0.22\textwidth}>{\raggedright\arraybackslash}p{0.14\textwidth}>{\raggedleft\arraybackslash}p{0.095\textwidth}>{\raggedleft\arraybackslash}p{0.105\textwidth}Y@{}}
\toprule
Case & City label & Genre family & First band & Emerg. & Geography note \\
\midrule
Birmingham heavy metal & Birmingham, United Kingdom & heavy metal & 1970 & 1977 & City-like label kept in baseline. \\
Tampa death metal & Tampa, United States & death metal & 1984 & 1987 & City-like label kept in baseline. \\
Bergen black metal & Bergen, Norway & black metal & 1991 & 1992 & City-like label kept in baseline. \\
Oslo black metal & Oslo, Norway & black metal & 1987 & 1990 & City-like label kept in baseline. \\
Gothenburg melodic death metal & Gothenburg, Sweden & melodic death metal & 1990 & 1996 & City-like label kept in baseline. \\
Bay Area thrash metal & Bay Area, United States & thrash metal & 1985 & 1992 & Region-like label excluded from city baseline. \\
\bottomrule
\end{tabularx}

\vspace{0.4em}
\parbox{0.96\textwidth}{\footnotesize Notes: These are operational scene dates from the project's
\texttt{city\_genre\_first\_appearance.csv} file, not claims about the exact historical birth date
of a genre. The point is narrower: the dates should be broadly plausible for canonical scenes, and
the geography note should clarify when a famous scene label is intentionally excluded because it is
a wider unit rather than a city-like observation.}
\end{table}


\subsection{Secondary digital-era checks}
\label{app:digital_checks}

The digital-era branch is kept as background only. A broad period split of the preferred
\texttt{roll-5} exact-year specification suggests that local spawning flow attenuates from
`0.0177` before `1995` to `0.0071` in `2005--2014`, and target-genre active bands attenuate from
`0.0977` to `0.0310`, while target-genre multi-band musicians stay stable at `0.0270`, `0.0312`,
and `0.0325`. That is a useful descriptive nuance, but not a different paper.

The broader internet and broadband probes are weaker. Country-level internet-adoption event studies
line up with higher later metal entry, but the effect is broad-based rather than cleanly
genre-specific, and the strongest black-metal case also shows a positive pretrend. The city-level
Ookla extension is narrower still: the current black-metal broadband probe covers only `47`
city-year observations across `33` cities in `2020--2022`, so it is too late and too small to
support a serious local internet mechanism claim. The practical implication is simple. These
checks are useful for scope discipline and for guarding against overclaiming, but they do not
strengthen the main scene result enough to belong in the core paper narrative.


\begin{thebibliography}{99}

\bibitem[Dixit and Stiglitz(1977)]{dixit_stiglitz_1977}
Dixit, Avinash K., and Joseph E. Stiglitz. 1977.
\newblock Monopolistic competition and optimum product diversity.
\newblock \emph{American Economic Review} 67(3): 297--308.

\bibitem[Duranton and Puga(2004)]{duranton_puga_2004}
Duranton, Gilles, and Diego Puga. 2004.
\newblock Micro-foundations of urban agglomeration economies.
\newblock In J. V. Henderson and J.-F. Thisse, eds., \emph{Handbook of Regional and Urban Economics}, vol. 4.

\bibitem[Ferreira and Waldfogel(2013)]{ferreira_waldfogel_2013}
Ferreira, Fernando, and Joel Waldfogel. 2013.
\newblock Pop internationalism: Has a half century of world music trade displaced local culture?
\newblock \emph{Economic Journal} 123(569): 634--664.

\bibitem[Franco and Filson(2006)]{franco_filson_2006}
Franco, April, and Darren Filson. 2006.
\newblock Spin-outs: Knowledge diffusion through employee mobility.
\newblock \emph{RAND Journal of Economics} 37(4): 841--860.

\bibitem[Glaeser et al.(1992)]{glaeser_etal_1992}
Glaeser, Edward L., Hedi D. Kallal, Jose A. Scheinkman, and Andrei Shleifer. 1992.
\newblock Growth in cities.
\newblock \emph{Journal of Political Economy} 100(6): 1126--1152.

\bibitem[Klement and Strambach(2019a)]{klement_strambach_eg_2019}
Klement, Benjamin, and Simone Strambach. 2019a.
\newblock Innovation in creative industries: Does (related) variety matter for the creativity of urban music scenes?
\newblock \emph{Economic Geography} 95(2): 146--167.

\bibitem[Klement and Strambach(2019b)]{klement_strambach_rs_2019}
Klement, Benjamin, and Simone Strambach. 2019b.
\newblock How do new music genres emerge? Diversification processes in symbolic knowledge bases.
\newblock \emph{Regional Studies} 53(10): 1447--1458.

\bibitem[Lena and Peterson(2008)]{lena_peterson_2008}
Lena, Jennifer C., and Richard A. Peterson. 2008.
\newblock Classification as culture: Types and trajectories of music genres.
\newblock \emph{American Sociological Review} 73(5): 697--718.

\bibitem[Moretti(2004)]{moretti_2004}
Moretti, Enrico. 2004.
\newblock Workers' education, spillovers, and productivity: Evidence from plant-level production functions.
\newblock \emph{American Economic Review} 94(3): 656--690.

\bibitem[Moen(2000)]{moen_2000}
Moen, Jarle. 2000.
\newblock Is mobility of technical personnel a source of R\&D spillovers?
\newblock \emph{NBER Working Paper} No. 7834.

\bibitem[Romer(1990)]{romer_1990}
Romer, Paul M. 1990.
\newblock Endogenous technological change.
\newblock \emph{Journal of Political Economy} 98(5, Part 2): S71--S102.

\bibitem[Scott(1999)]{scott_1999}
Scott, Allen J. 1999.
\newblock The US recorded music industry: On the relations between organization, location, and creativity in the cultural economy.
\newblock \emph{Environment and Planning A} 31(11): 1965--1984.

\bibitem[Storper and Venables(2004)]{storper_venables_2004}
Storper, Michael, and Anthony J. Venables. 2004.
\newblock Buzz: Face-to-face contact and the urban economy.
\newblock \emph{Journal of Economic Geography} 4(4): 351--370.

\bibitem[Waldfogel(2015)]{waldfogel_2015}
Waldfogel, Joel. 2015.
\newblock Digitization and the quality of new media products: The case of music.
\newblock In A. Goldfarb, S. Greenstein, and C. Tucker, eds., \emph{Economic Analysis of the Digital Economy}.

\bibitem[Weitzman(1998)]{weitzman_1998}
Weitzman, Martin L. 1998.
\newblock Recombinant growth.
\newblock \emph{Quarterly Journal of Economics} 113(2): 331--360.

\end{thebibliography}


\end{document}
